\section{For Discussion}

A working framework and a settled conclusion are different things. The five points below bear most directly on whether a specific deployment decision is defensible right now; the complete set of fourteen --- including methodological questions more relevant to refining the framework than to a specific decision --- is in Appendix B.

\textbf{2. How much practical weight a validated confidence signal should actually get.} The position in Section 9 --- use confidence for triage only where independently validated for the specific model, task, and deployment --- is defensible as stated, but it's fair to ask how often that validation work will actually get done in practice. One view: this caveat is close to vacuous in most real deployments, since the validation work is expensive, ongoing, and rarely prioritized, so the practical guidance collapses to "check everything" regardless of the theoretical allowance. The opposing view: as calibration research matures and deployment-specific validation becomes cheaper and more standard, this caveat will do real work, and building the \emph{option} into the framework now is worth the complexity even if few deployments currently exercise it.

\textbf{4. How strong the AISI/OpenAI evidence base actually generalizes.} This report treats the combination of the OpenAI incident and the AISI cross-lab finding as materially stronger evidence than a single incident for the claim that frontier models commonly route around inconvenient evaluation boundaries. That's a defensible strengthening --- but "cheating on a cybersecurity capability evaluation, in a setting with deliberately relaxed safety controls" is still a specific context, and how far this generalizes to, for example, an agent operating in a normal commercial deployment with standard (not relaxed) safety measures and a non-adversarial task is genuinely open. The AISI and OpenAI findings support taking containment seriously as a default posture; they do not by themselves establish the base rate of boundary-defeat behavior in ordinary, non-adversarial-benchmark conditions, and that base rate matters a great deal for calibrating how much containment engineering effort is proportionate for a specific low-stakes deployment versus a specific high-stakes one.

\textbf{8. Whether the six-property verification schema and the sixteen-field worksheet risk manufacturing the same false precision they're meant to cure.} Section 13's schema and Section 14's worksheet are genuinely useful as prompts for what to ask. But a filled-in cell reading "sensitivity: 80\%" that was actually a reviewer's on-the-spot guess, never measured, is arguably worse than an honest blank --- it carries the visual authority of a measured figure without the substance of one, which is exactly the confidence-versus-reliability gap Section 9 spends considerable effort warning about. Section 14's proportionality caveat is this report's attempt to guard against that outcome, but whether that caveat will actually be respected in practice, under real organizational pressure to show a completed-looking assessment, is an open question about implementation discipline that no amount of caveat-writing can settle in advance.

\textbf{11. Whether a genuinely "strong human baseline" can actually be established for most of the tasks this framework cares about.} Section 2's comparative standard depends on knowing what a well-rested, properly trained human actually achieves on a given task --- but for a great deal of professional-judgment work, that figure has never been rigorously measured at all. Organizations rarely track their own error rates, omission rates, or consistency-across-equivalent-cases with anything like the discipline Section 2 asks of the AI side of the comparison. There's a real risk that the comparison ends up asymmetric in practice: the AI system gets measured carefully because it's new and being scrutinized, while the human baseline it's compared against is estimated, assumed, or drawn from whatever informal impression happens to be on hand. Where that happens, a "beats the human baseline" conclusion may be true only because the baseline was never rigorously established in the first place, not because the comparison was actually won. This is a caution about implementation, not a flaw in the standard itself.

\textbf{12. Whether comparative framing, despite its own stated safeguards, is more resistant to being gamed than a stricter standard would be.} Section 2 explicitly warns against measuring an AI system against a weak human comparator, and explicitly requires tail-risk constraints alongside average performance specifically to prevent a good-on-average system with an intolerable tail from being waved through. Those safeguards are well-designed on paper. Whether they hold up under real deployment pressure --- where there's an organizational incentive to declare a comparison won, a baseline is inconvenient or expensive to measure rigorously, and tail risks are by definition rare enough to not yet have shown up in whatever evidence exists at decision time --- is a separate question the framework's own logic can't settle.
